\section{STL esencial}

\begin{center}
\begin{tabular}{lll}
Estructura & Operaciones principales & Complejidad \\
\hline
\texttt{vector} & acceso / push al final & $O(1)$ / amortizado $O(1)$ \\
\texttt{deque} & push/pop en ambos extremos & $O(1)$ \\
\texttt{stack, queue} & push/pop/top/front & $O(1)$ \\
\texttt{priority\_queue} & push/pop/top & $O(\log n)$ / $O(1)$ \\
\texttt{set, map} & find/insert/erase por clave & $O(\log n)$ \\
\texttt{unordered\_*} & find/insert/erase & promedio $O(1)$ \\
\end{tabular}
\end{center}

\begin{lstlisting}[style=cpp]
void stl_examples(vector<ll> &v, ll x) {
    sort(v.begin(), v.end());
    v.erase(unique(v.begin(), v.end()), v.end());

    auto lo = lower_bound(v.begin(), v.end(), x); // primero >= x
    auto hi = upper_bound(v.begin(), v.end(), x); // primero > x
    if (lo != v.end()) { ll value = *lo; (void)value; }
    (void)hi;

    priority_queue<int> max_heap;
    priority_queue<int, vector<int>, greater<int>> min_heap;

    ll total = accumulate(v.begin(), v.end(), 0LL);
    (void)max_heap; (void)min_heap; (void)total;
}
\end{lstlisting}

\warning{\texttt{reserve(n)} reserva capacidad pero no crea elementos. Un
arreglo local suele vivir en el stack; uno global o \texttt{static}, no.}
